\begin{definition}[The network adversary]
  \label{def:aba-netadv}
  \lean{PLTS.ABA.Net.NetState}\lean{PLTS.ABA.Net.NetState.corrupt}\lean{PLTS.ABA.Net.NetStep}\lean{PLTS.ABA.Net.netAdv}\leanok
  \uses{def:aba-procnet, def:aba-labels, def:aba-params}
  The network adversary $\mathrm{netAdv}(P)$ holds what a process may not see: the round-tagged sent sets $sent(r,j)$ of the stage messages $j$ has multicast in round $r$ (D5), the sent sets $dpool(j)$ of the DECIDED payloads $j$ has multicast (D12$'$), and the corrupted set $F$ with its budget, the guard $k \notin F \wedge |F| < f$ on the $\mathrm{fail}$ row (D1).
  It supplies the other half of every send and every delivery, the delivery guards $m \in sent(r,j)$ and $b \in dpool(j)$ being the soundness of the authenticated network; it carries the authorisation $k \in F$ of each of the five Byzantine handshake rows (D11), and the two injections $\mathrm{byzG}$ and $\mathrm{byzD}$, by which a corrupted sender sent sets an arbitrary stage message or either DECIDED bit and so equivocates in those sent sets.
  Its row $\mathrm{retByz}$ carries the same authorisation and lets a corrupted process return either bit with none of the DECIDED evidence the honest return row asks for; the replaced program's self-loop is the other half of it (D23).
  Full rule table: \leandecl{PLTS.ABA.Net.netAdv}.
\end{definition}
